% =============================================================================
% UQFF v5.78  Section Template :: T-Lambda
% Cosmological-constant / dark-energy / vacuum-density anchored sections
% -----------------------------------------------------------------------------
% USAGE
%   1. Copy this file to whitepapers/PAPER_<NNNN>_<slug>.tex
%   2. Replace every <PLACEHOLDER> token (search for the literal '<')
%   3. Keep every CANONICAL :: line VERBATIM (the closure parser keys on them)
%   4. Compile with pdflatex direct (NO PANDOC).  Two passes.
% -----------------------------------------------------------------------------
% CANONICAL CONSTANTS (do not edit; mirrored from dpm_vacuum_manifold.py v3.0)
%   RHO_VAC_SCM = 4*sqrt(pi)*1e-37 = 7.0898154036e-37 J/m^3   (G9 structural)
%   RHO_VAC_UA  = 10 * RHO_VAC_SCM = 7.0898154036e-36 J/m^3   (G7, |SO(5)|)
%   BETA_I = 0.60     KAPPA = 5e-4 / day     [SSq] = 0.57
%   H_SCm  = 0.99     U_UA  = 1.0e-4         F_TRZ = 0.1
%   k_eta  = F_TRZ^113 = 1.0e-113
% =============================================================================
\documentclass[11pt]{article}
\usepackage[margin=1in]{geometry}
\usepackage{amsmath,amssymb,amsthm}
\usepackage{booktabs,longtable,array}
\usepackage{listings,xcolor}
\usepackage[colorlinks=true,linkcolor=blue,urlcolor=blue,citecolor=blue]{hyperref}
\lstset{basicstyle=\ttfamily\small,breaklines=true,frame=single,
        language=Python,showstringspaces=false,columns=fullflexible}

\title{<PAPER TITLE>: \\ Cosmological-Constant / Vacuum-Density Closure under v5.78}
\author{Star-Magic UQFF \\ PAPER\_<NNNN> \,/\, Session <SESSION>}
\date{<DATE>}

% --- Closure-tracker banner (parser keys on these lines, do not edit format) --
% CLOSURE :: <closure_label> :: predicted=<X> observed=<Y> error_pct=<Z>
% TEMPLATE :: T-Lambda
% CANONICAL :: rho_vac_scm = 7.0898154036e-37 J/m^3
% CANONICAL :: rho_vac_ua  = 7.0898154036e-36 J/m^3

\begin{document}
\maketitle

\begin{abstract}
<2--6 sentence abstract.  State: (a) the cosmological-constant/dark-energy
anchor being closed, (b) the canonical UQFF expression involving
$\rho_{\mathrm{vac,SCm}}$ or $\rho_{\mathrm{vac,UA}}$, (c) the observational
target (Planck 2018, DES, eBOSS, JWST, etc.), and (d) the resulting
predicted-vs-observed error.>
\end{abstract}

\section{Anchor and Target}
\textbf{Observational anchor.} <e.g. Planck 2018 $\Omega_\Lambda h^2 = \ldots$ at $z=0$.>\\
\textbf{Reference.} <DOI / arXiv / catalog ID> --- this MUST be a real citation.\\
\textbf{UQFF target quantity.} <name and symbol>.

\section{Canonical Equation}
\begin{equation}
  \boxed{\;<UQFF EQUATION involving \rho_{\mathrm{vac,SCm}} or \rho_{\mathrm{vac,UA}}>\;}
  \label{eq:tlambda-main}
\end{equation}
The structural derivation enters via $\rho_{\mathrm{vac,SCm}} = 4\sqrt{\pi}\cdot 10^{-37}\,\mathrm{J/m^3}$
under the G9 structural closure (PAPER\_409, hydrogen geometry); the
$|SO(5)|=10$ ratio to $\rho_{\mathrm{vac,UA}}$ is G7 (PAPER\_1160).

\section{Numerical Closure}
\begin{longtable}{lll}
\toprule
Quantity & Predicted (UQFF v5.78) & Observed \\
\midrule
<quantity 1> & <value> & <value $\pm$ error> \\
<quantity 2> & <value> & <value $\pm$ error> \\
\bottomrule
\end{longtable}

\paragraph{Error.} $\varepsilon = |P-O|/O = <Z>\%$.

\section{Cross-references}
\begin{itemize}
  \item Canonical vacuum source: \texttt{dpm\_vacuum\_manifold.py} L97--98 (v3.0).
  \item Calculator: <e.g. \texttt{CondensedPhysics3.py::CosmologicalConstantCalculator}>
  \item Prior closure(s): PAPER\_<NN>, PAPER\_<MM>.
\end{itemize}

\section{Reproducibility}
\begin{lstlisting}
from dpm_vacuum_manifold import RHO_VAC_SCM, RHO_VAC_UA
# show the one-liner that reproduces the predicted value
\end{lstlisting}

\bibliographystyle{plain}
\begin{thebibliography}{9}
\bibitem{planck2018} Planck Collaboration, \emph{Planck 2018 results VI. Cosmological parameters}, A\&A 641, A6 (2020). \texttt{arXiv:1807.06209}.
\bibitem{paper409}   PAPER\_409 (Star-Magic): \emph{Hydrogen-Geometry Derivation of $\rho_{\mathrm{vac,SCm}}$}.
\bibitem{paper1160}  PAPER\_1160 (Star-Magic): \emph{F\_TRZ G7 Structural Closure / |SO(5)| Ratio}.
\end{thebibliography}

\end{document}
